\section{Classical reductions with one index equal to two}
\label{recip:one2:sec:general}

This section settles the three Gaussian weight-four triple identities
printed in Chapter~4 of the source manuscript.  The proof gives a uniform
formula for every multiple polylogarithm whose indices consist of one
$2$ and any number of $1$'s.  The general phenomenon is classical:
Borwein and Straub give the complementary-argument formula for
$\Li_{2,\{1\}^{p-1}}$ in \cite[Eq.~(33)]{gaussian:BorweinStraub2015}, including
the weight-four instance in their Eq.~(34).  The results below reorganize
that reduction at the reciprocal complementary argument, make the
dependence on the position of the $2$ explicit, and prove the specific
Gaussian formulas in the manuscript.  We do not claim that classical
reducibility of this family is a new phenomenon.

Throughout this section, a single displayed argument abbreviates the
remaining arguments $1$:
\[
 \Li_{s_1,\ldots,s_d}(z)
 :=\Li_{s_1,\ldots,s_d}(z,1,\ldots,1).
\]
The defining series has strictly decreasing summation indices.  Set
\[
 U_{a,b}(z)=\Li_{\{1\}^{a},2,\{1\}^{b}}(z),\qquad
 H_p(z)=\Li_{2,\{1\}^{p-1}}(z),
\]
where $a,b\ge0$ and $p\ge1$.  The weight of $U_{a,b}$ is
$N=a+b+2$, and its series depth is $N-1$.

\subsection{A logarithmic moment and its branch conventions}

We use principal branches and first work on
\[
 \mathcal D=\mathbb C\setminus[0,\infty),\qquad
 w=\frac1{1-z},\qquad L=\Log w=-\Log(1-z).
\]
On this domain, $w$ avoids the standard polylogarithm cut $[1,\infty)$
and the logarithm cut $(-\infty,0]$.  All terms in the formulas below
are therefore single-valued analytic functions.  The formulas also
extend to the usual slit plane by consistent continuation; no such
extension is needed at $z=i$ or at a nonreal root of unity.

\begin{lemma}[Logarithmic moment]
\label{recip:one2:lem:moment}
Let $q(t)=-\Log(1-t)$.  For an integration path from $0$ to $z$ avoiding
$\{0,1\}$ apart from its initial point $0$, with the branches
continued along that path,
\begin{equation}
 U_{a,b}(z)=\frac1{a!(b+1)!}
 \int_0^z\frac{(q(z)-q(t))^a q(t)^{b+1}}{t}\,dt.
 \label{recip:one2:eq:moment}
\end{equation}
Consequently,
\begin{equation}
 U_{a,b}(z)=\sum_{j=0}^a(-1)^j
 \binom{b+j+1}{j}\frac{q(z)^{a-j}}{(a-j)!}\,H_{b+j+1}(z).
 \label{recip:one2:eq:triangular}
\end{equation}
\end{lemma}
\begin{proof}
The iterated-integral word of $U_{a,b}$ is
$1^a0\,1^{b+1}$, with forms $dt/(1-t)$ and $dt/t$ attached to $1$ and
$0$, respectively.  The innermost $b+1$ identical forms contribute
$q(t)^{b+1}/(b+1)!$.  The $a$ identical forms between $t$ and $z$
contribute $(q(z)-q(t))^a/a!$.  Integrating the single remaining $0$
form proves \eqref{recip:one2:eq:moment}.  Equivalently, one may verify the
same formula by differentiating successively with respect to $z$ and
using the zero initial values at $z=0$.  Expanding the binomial and
using
\[
 H_p(z)=\frac1{p!}\int_0^z\frac{q(t)^p}{t}\,dt
\]
proves \eqref{recip:one2:eq:triangular}.
\end{proof}

\begin{theorem}[Explicit reduction for every position of the $2$]
\label{recip:one2:thm:closed}
For $a,b\ge0$, $N=a+b+2$, and $z\in\mathcal D$,
\begin{align}
 U_{a,b}(z)
 &=\sum_{k=a+1}^{N}(-1)^{a+k}\binom{k-1}{a}
       \frac{L^{N-k}}{(N-k)!}\Li_k(w)
       -\frac{L^N}{N!}\notag\\
 &\quad+(-1)^{b+1}\sum_{j=0}^{a}\binom{b+j+1}{j}
       \frac{L^{a-j}}{(a-j)!}\zeta(b+j+2).
 \label{recip:one2:eq:closed}
\end{align}
In particular,
\begin{equation}
 H_p(z)=(-1)^p\zeta(p+1)-\frac{L^{p+1}}{(p+1)!}
 +\sum_{k=1}^{p+1}(-1)^k
       \frac{L^{p+1-k}}{(p+1-k)!}\Li_k(w).
 \label{recip:one2:eq:height}
\end{equation}
\end{theorem}
\begin{proof}
We first prove \eqref{recip:one2:eq:height} for real $z<0$, so that
$0<w<1$ and every intermediate integral is real.  Substitute
$v=1/(1-t)$ in the logarithmic moment.  Then
\[
 \frac{dt}{t}=\frac{dv}{v(v-1)},\qquad q(t)=\log v,
\]
and
\[
 H_p(z)=\frac1{p!}\int_1^w\frac{\log^p v}{v(v-1)}\,dv.
\]
The differential identity
\begin{equation}
 \frac{d}{dv}\left[
 \sum_{k=1}^{p+1}\frac{(-1)^{k-1}p!}{(p+1-k)!}
       \log^{p+1-k}v\,\Li_k(v)\right]
       =\frac{\log^p v}{1-v}
 \label{recip:one2:eq:primitive}
\end{equation}
follows by cancellation of adjacent terms, using
$\Li_1(v)=-\log(1-v)$ and $\Li'_k(v)=\Li_{k-1}(v)/v$ for $k\ge2$.
Since
$1/[v(v-1)]=-1/v-1/(1-v)$, integration gives
\eqref{recip:one2:eq:height}.  At $v=1$, every term in
\eqref{recip:one2:eq:primitive} with a positive power of $\log v$ tends to
zero.  This includes the apparently singular term
$\log^p v\,\Li_1(v)$, which is
$O((v-1)^p\log(1-v))$ and tends to zero because $p\ge1$.
The sole nonzero endpoint term is $(-1)^pp!\zeta(p+1)$.
Thus the constant in \eqref{recip:one2:eq:height} is fixed, not inferred
from a differential identity alone.  Analytic continuation on the
connected domain $\mathcal D$ proves that identity throughout the
stated domain.

Insert \eqref{recip:one2:eq:height} in
\eqref{recip:one2:eq:triangular}, using $q(z)=L$.  The zeta terms immediately
give the second line of \eqref{recip:one2:eq:closed}.  The coefficient of
$(-1)^k L^{N-k}\Li_k(w)$ is
\[
 \sum_{j=0}^a\frac{(-1)^j\binom{b+j+1}{j}}
                    {(a-j)!(b+j+2-k)!}
 =\frac{(-1)^a\binom{k-1}{a}}{(N-k)!}.
\]
Terms containing a factorial of a negative integer are omitted.  To
verify this finite identity, write its left side as
\[
 \frac{(k-1)!}{(b+1)!a!}
 \sum_{j=0}^a(-1)^j\binom aj\binom{b+j+1}{k-1}
\]
and extract the coefficient of $x^{k-1}$ in
\[
 (1+x)^{b+1}\sum_{j=0}^a(-1)^j\binom aj(1+x)^j
   =(-x)^a(1+x)^{b+1}.
\]
The coefficients vanish when $k\le a$, explaining the lower endpoint
$a+1$ in the sum.  Finally, the coefficient of $-L^N$ equals
\begin{align*}
 \sum_{j=0}^a\frac{(-1)^j\binom{b+j+1}{j}}
                    {(a-j)!(b+j+2)!}
 &=\frac1{a!(b+1)!}\int_0^1t^{b+1}(1-t)^a\,dt\\
 &=\frac1{N!}.
\end{align*}
This proves the complete formula.
\end{proof}

For the Gaussian endpoint the path itself can be inspected explicitly:
\[
 z=ix\quad(0\le x\le1),\qquad
 w(x)=\frac{1+ix}{1+x^2}.
\]
For $x>0$ this path lies in the upper half-plane and does not meet the
polylogarithm cut.  Its initial point is $w(0)=1$, where the endpoint
limits in the proof apply, and its final point is $(1+i)/2$.  This
verifies the branch choices directly, without using an inversion
formula with an unspecified logarithmic correction.

The boundary series at $z=i$ agree with these analytic values.
Indeed the finite nested harmonic sums in a coefficient satisfy a
bound $O((1+\log n)^{d-1})$.  After division by $n^{s_1}$ their
coefficients tend to zero and have summable successive differences.
Summation by parts against the bounded partial sums of $i^n$ proves
convergence, even when $s_1=1$; Abel's theorem then identifies the
series with the radial analytic limit.

\begin{corollary}[A short formula when the $2$ is last]
\label{recip:one2:cor:last}
For every $N\ge2$,
\begin{align}
 U_{N-2,0}(z)
 &=(N-1)\Li_N(w)-L\Li_{N-1}(w)-\frac{L^N}{N!}\notag\\
 &\quad-\sum_{m=0}^{N-2}(N-1-m)\zeta(N-m)\frac{L^m}{m!}.
 \label{recip:one2:eq:last}
\end{align}
Thus only two nonconstant classical polylogarithms are needed in this
subfamily, regardless of series depth.
\end{corollary}
\begin{proof}
Set $a=N-2$, $b=0$ in \eqref{recip:one2:eq:closed} and reindex the zeta sum.
\end{proof}

\begin{corollary}[The top two series-depth layers]
\label{recip:one2:cor:toplayers}
Every single-endpoint multiple polylogarithm of weight $N\ge2$ and
series depth at least $N-1$ has an explicit expression in classical
polylogarithms at $(1-z)^{-1}$, zeta values, and logarithms, on the
domain $\mathcal D$.
\end{corollary}
\begin{proof}
Positive indices with sum $N$ and length $N$ are all $1$, and
$\Li_{\{1\}^N}(z)=L^N/N!$. Those with length $N-1$ have exactly one
index $2$ and all other indices $1$, so Theorem~\ref{recip:one2:thm:closed}
applies. These exhaust the possibilities.
\end{proof}

This statement concerns the literal depth of the defining series.
It does not identify that depth with a motivic filtration or prove a
minimal depth for the resulting numerical constants.

\subsection{One possible new direction at each weight}

Put
\[
 \ell=\log2,\qquad w=\frac{1+i}{2},\qquad
 \mu_k=\Re\Li_k(w),\qquad\lambda_k=\Im\Li_k(w).
\]
For $N\ge2$, let $\mathcal P_N$ be the rational span of products of
at least two positive-weight atoms selected from
\[
 \pi,\ell,\quad \zeta(k),\mu_k,\lambda_k\quad(2\le k<N),
\]
whose assigned weights sum to $N$; the assigned weights are $1$ for
$\pi,\ell$ and $k$ for the atoms with subscript or argument $k$.
This definition is a concrete span of numerical products.  It does
not assume that numerical periods form a direct sum by weight.

\begin{corollary}[Signed binomial coefficients modulo products]
\label{recip:one2:cor:primitive}
For $a+b+2=N$, set
$c_{N,a}=(-1)^{N+a}\binom{N-1}{a}$.  Then
\begin{align}
 U_{a,b}(i)&\equiv c_{N,a}\bigl(\Li_N(w)-\zeta(N)\bigr)
        \pmod{\mathcal P_N+i\mathcal P_N},\label{recip:one2:eq:complexprimitive}\\
 \Re U_{a,b}(i)&\equiv c_{N,a}\bigl(\mu_N-\zeta(N)\bigr)
        \pmod{\mathcal P_N},\label{recip:one2:eq:realprimitive}\\
 \Im U_{a,b}(i)&\equiv c_{N,a}\lambda_N
        \pmod{\mathcal P_N}.\label{recip:one2:eq:imagprimitive}
\end{align}
In particular, the imaginary parts of this entire family contribute
at most one direction after products of the stated lower-weight
atoms have been factored out.
\end{corollary}
\begin{proof}
At $z=i$, $L=-\ell/2+i\pi/4$, and
$\Li_1(w)=-\overline L$.  In \eqref{recip:one2:eq:closed}, every term is a
product of lower-weight atoms except the term $k=N$ and the term
$j=a$ in the zeta sum.  Their respective coefficients are
$c_{N,a}$ and $-c_{N,a}$.  Taking real and imaginary parts proves the
statements.
\end{proof}

The coefficients of the highest-weight half-point imaginary value are
\[
 \begin{array}{c|l}
 N&c_{N,0},\ldots,c_{N,N-2}\\\hline
 4&1,-3,3\\
 5&-1,4,-6,4\\
 6&1,-5,10,-10,5.
 \end{array}
\]
The corollary is an upper bound, and proves neither that $\lambda_N$
is nonzero in the quotient nor that it is independent of constants
from other argument families.  It also does not identify series
depth with an additive motivic depth filtration.

\subsection{A sixth-root specialization in every position}
\label{recip:one2:sec:sixthroot}

The special role of sixth roots in relations for Nielsen polylogarithms
is studied by Borwein and Straub~\cite{gaussian:BorweinStraub2015}.  The preceding
position formula gives a direct specialization in that established
setting, including an explicit rational coefficient for one component
at every weight.  Throughout this subsection,
\[
 \omega=e^{i\pi/3},\qquad N=a+b+2,\qquad a,b\ge0,
\]
and all logarithms and polylogarithms use their principal branches.

\begin{corollary}[Sixth-root closure and a pure parity component]
\label{recip:one2:cor:sixthroot}
At $z=\omega$, formula~\eqref{recip:one2:eq:closed} reduces
$U_{a,b}(\omega)$ to classical values $\Li_k(\omega)$, zeta values,
and powers of $i\pi$, since
\begin{equation}
 \frac1{1-\omega}=\omega,\qquad
 -\Log(1-\omega)=\Log\omega=\frac{i\pi}{3}.
 \label{recip:one2:eq:sixthfixed}
\end{equation}
Moreover,
\begin{equation}
 U_{a,b}(\omega)+(-1)^N\overline{U_{a,b}(\omega)}
       =R_{N,a}(i\pi)^N,\qquad R_{N,a}\in\mathbb Q.
 \label{recip:one2:eq:sixthparity}
\end{equation}
Thus $\Re U_{a,b}(\omega)$ is a rational multiple of $\pi^N$
when $N$ is even, and $\Im U_{a,b}(\omega)$ is a rational multiple
of $\pi^N$ when $N$ is odd.  Explicitly, with Bernoulli polynomials
defined by
\[
 \frac{t e^{xt}}{e^t-1}
    =\sum_{k\ge0}B_k(x)\frac{t^k}{k!},\qquad B_k=B_k(0),
\]
one may take
\begin{align}
 R_{N,a}
 &=-\frac{2}{3^N N!}
 -\sum_{k=a+1}^{N}
   \frac{(-1)^{a+k}\binom{k-1}{a}\,2^k B_k(1/6)}
        {3^{N-k}(N-k)!k!}\notag\\
 &\quad+(-1)^{N-a}
   \sum_{\substack{N-a\le m\le N\\m\ \mathrm{even}}}
   \frac{\binom{m-1}{N-a-1}\,2^m B_m}
        {3^{N-m}(N-m)!m!}.
 \label{recip:one2:eq:sixthcoefficient}
\end{align}
The same formulas at $\overline\omega$ follow by complex conjugation.
\end{corollary}
\begin{proof}
The identities in \eqref{recip:one2:eq:sixthfixed} follow from
$1-\omega=\overline\omega$.  Along the radial segment
$z=r\omega$, $0\le r\le1$, the real part of $1-z$ is at least
$1/2$, so the indicated logarithm has no branch crossing.
The transformed path lies in the upper half-plane for $r>0$,
starts at $1$, and ends at $\omega$; the endpoint limits already
used to prove \eqref{recip:one2:eq:closed} apply.  Consequently that
formula has exactly the stated principal values.  The ordinary
boundary multiple series agree with their radial limits by the
same summation-by-parts argument used at $z=i$.

Put $T=i\pi$ and $L=T/3$.  The classical Bernoulli Fourier identity is
\begin{equation}
 \Li_k(e^{2\pi i x})+(-1)^k\Li_k(e^{-2\pi i x})
       =-\frac{(2\pi i)^k}{k!}B_k(x),
 \qquad 0<x<1,\quad k\ge1.
 \label{recip:one2:eq:bernoullifourier}
\end{equation}
For $k=1$, this follows directly from $\Li_1(z)=-\Log(1-z)$
on the stated arc.  Successive integration, with the zero average
of both sides fixing the constants, proves every higher case.
Since $\overline L=-L$, the parity projection of a classical term
in \eqref{recip:one2:eq:closed} is
\begin{align*}
 &L^{N-k}\Li_k(\omega)
       +(-1)^N\overline{L^{N-k}\Li_k(\omega)}\\
 &\hspace{1cm}=L^{N-k}
    \bigl(\Li_k(\omega)+(-1)^k\Li_k(\overline\omega)\bigr)
   =-\frac{2^k B_k(1/6)}{3^{N-k}k!}\,T^N.
\end{align*}
The polynomial term $-L^N/N!$ contributes
$-2T^N/(3^N N!)$.  A zeta term with zeta weight $m$ contributes
zero when $m$ is odd.  For even $m\ge2$, use
\[
 \frac{2\zeta(m)}{(i\pi)^m}=-\frac{2^m B_m}{m!}.
\]
In the last sum of \eqref{recip:one2:eq:closed}, write
$m=b+j+2=N-a+j$.  Its binomial coefficient becomes
$\binom{m-1}{N-a-1}$, while its sign after the even-zeta evaluation
is $(-1)^b=(-1)^{N-a}$.  Combining these three contributions proves
\eqref{recip:one2:eq:sixthcoefficient} and hence
\eqref{recip:one2:eq:sixthparity}.
\end{proof}

There is a precise geometric reason that this specialization closes
on a unit-circle argument.  If $|z|=1$ and $z\ne1$, then
\[
 \left|\frac1{1-z}\right|=1
 \quad\Longleftrightarrow\quad
 |1-z|^2=2-2\Re z=1
 \quad\Longleftrightarrow\quad z\in\{\omega,\overline\omega\}.
\]
Both of these points are fixed by $z\mapsto(1-z)^{-1}$.
This classifies the unit-circle endpoint geometry of this particular
transformation.

\begin{corollary}[The two weight-three values]
\label{recip:one2:ex:sixththree}
Define the ordinary Clausen value by its absolutely convergent series
\[
 C=\operatorname{Cl}_2(\pi/3)
   =\sum_{n\ge1}\frac{\sin(n\pi/3)}{n^2}.
\]
Then
\begin{align}
 U_{0,1}(\omega)=\Li_{2,1}(\omega)
  &=\frac23\zeta(3)-\frac\pi3 C+\frac{i\pi^3}{324},
    \label{recip:one2:eq:sixth21}\\
 U_{1,0}(\omega)=\Li_{1,2}(\omega)
  &=-\frac43\zeta(3)+\frac\pi3 C+\frac{i\pi^3}{324}.
    \label{recip:one2:eq:sixth12}
\end{align}
\end{corollary}
\begin{proof}
The classical values needed are
\[
 \Li_1(\omega)=\frac{i\pi}{3},\qquad
 \Li_2(\omega)=\frac{\pi^2}{36}+iC,\qquad
 \Li_3(\omega)=\frac{\zeta(3)}3+\frac{5i\pi^3}{162}.
\]
The real dilogarithm component and imaginary trilogarithm component
follow from \eqref{recip:one2:eq:bernoullifourier}.  The real
trilogarithm component follows by grouping the defining series
into residue classes, or from the distribution identity
\[
 \Li_s(\omega)+\Li_s(\overline\omega)
   =(1-2^{1-s})(1-3^{1-s})\zeta(s),\qquad s>1.
\]
For example, summing over the three cube roots gives
$(3^{1-s}-1)\zeta(s)$ for the two nontrivial cube roots;
applying the two-term distribution formula to the pair of sixth
roots multiplies that sum by $2^{1-s}-1$.
Now \eqref{recip:one2:eq:height} at $p=2$ gives, with $L=i\pi/3$,
\[
 U_{0,1}(\omega)
   =\zeta(3)-\Li_3(\omega)+L\Li_2(\omega)-\frac23L^3.
\]
Substitution yields \eqref{recip:one2:eq:sixth21}.  The shuffle identity
\[
 \Li_1(\omega)\Li_2(\omega)
     =U_{1,0}(\omega)+2U_{0,1}(\omega)
\]
then proves \eqref{recip:one2:eq:sixth12}.  In both positions,
\eqref{recip:one2:eq:sixthcoefficient} gives $R_{3,a}=-1/162$.
\end{proof}

The independent script \texttt{code/independent/one\_two/sixth\_root.py}
computes \eqref{recip:one2:eq:sixthcoefficient} with exact rational
arithmetic and compares the specialized closed formula, its pure
parity component, and the displayed weight-three examples with the
logarithmic moment.  It covers every position of the $2$ at weights
$2$ through $8$, at both $\omega$ and $\overline\omega$, using
$100$ working decimal digits.  Across these $56$ cases, the largest
closed-form versus moment residual is below $6.59\cdot10^{-100}$;
the largest parity-projection residual is below $7.15\cdot10^{-102}$.
The four checks of the two explicit examples and their conjugates
have residual below $2.16\cdot10^{-101}$.  These are numerical checks
of the proved identities, rather than rigorous quadrature error bounds.


\subsection{Further exact identities at weights five and six}

The same calculation gives formulas beyond the three displayed
manuscript rows.  They involve the unavoidable output of this
particular reduction, namely higher classical values at $w$; no
independence claim about those values is implied.

\begin{corollary}[Two higher-weight Gaussian reductions]
\label{recip:one2:cor:56}
With $H_4=\Li_{2,1,1,1}$ and $H_5=\Li_{2,1,1,1,1}$,
\begin{align}
 \Im H_4(i)
 &=-\lambda_5-\frac{\ell\lambda_4}{2}+\frac{\pi\mu_4}{4}
       +\frac{\pi^2-4\ell^2}{32}\lambda_3\notag\\
 &\quad+\frac{G\ell(3\pi^2-4\ell^2)}{192}
       +\frac{35\pi\ell\zeta(3)}{512}
       +\frac{\pi\ell^2(\ell^2-\pi^2)}{384}
       -\frac{17\pi^5}{92160},
       \label{recip:one2:eq:H4}\\
 \Im H_5(i)
 &=\lambda_6+\frac{\ell\lambda_5}{2}-\frac{\pi\mu_5}{4}
       +\frac{4\ell^2-\pi^2}{32}\lambda_4
       -\frac{\pi\ell\mu_4}{8}\notag\\
 &\quad+\frac{\ell(4\ell^2-3\pi^2)}{192}\lambda_3
       +\frac{G(16\ell^4-24\ell^2\pi^2+\pi^4)}{6144}\notag\\
 &\quad+\frac{35\pi(\pi^2-12\ell^2)\zeta(3)}{24576}
       +\frac{\pi\ell^3(5\pi^2-3\ell^2)}{5760}.
       \label{recip:one2:eq:H5}
\end{align}
\end{corollary}
\begin{proof}
Use $p=4$ and $p=5$ in \eqref{recip:one2:eq:height}, substitute
$L=-\ell/2+i\pi/4$ and Lemma~\eqref{gaussian:eq:halfpoint-auxiliary}, and take
imaginary parts.  The accompanying coefficient generator performs
these polynomial operations over $\mathbb Q$ and also emits every
position of the $2$ at weights four, five, and six.
\end{proof}

The real parts can be treated by the same formula.  For example,
\begin{equation}
 \Re\Li_{2,1,1}(i)=\mu_4+\frac{\pi\lambda_3}{4}
 +\frac{\pi G\ell}{8}+\frac{35\ell\zeta(3)}{128}
 +\frac{\ell^4}{384}-\frac{11\pi^2\ell^2}{768}
 -\frac{1249\pi^4}{92160}.
 \label{recip:one2:eq:real211}
\end{equation}
The machine-readable tables include both real and imaginary parts,
so this companion family can be reproduced without manual
transcription.
